\documentclass[a4paper,12pt]{article}

\usepackage{cmap}		
\usepackage[utf8]{inputenc}			
\usepackage[english,russian]{babel}
\usepackage{framed}
\usepackage{hyperref}
\usepackage{amsmath}
\usepackage{graphicx}
\usepackage[colorinlistoftodos]{todonotes}
\usepackage{wrapfig}
\usepackage{lipsum}
\usepackage{listings}
\usepackage{color}
\usepackage{indentfirst}
\usepackage{times}
\usepackage{textcomp}


\definecolor{mygray}{rgb}{0.4,0.4,0.4}
\definecolor{mygreen}{rgb}{0,0.8,0.6}
\definecolor{myorange}{rgb}{1.0,0.4,0}

\lstdefinestyle{customc}{
  belowcaptionskip=1\baselineskip,
  breaklines=true,
  frame=L,
  xleftmargin=\parindent,
  language=C,
  showstringspaces=false,
  basicstyle=\footnotesize\ttfamily,
  keywordstyle=\bfseries\color{green!40!black},
  commentstyle=\itshape\color{purple!40!black},
  identifierstyle=\color{blue},
  stringstyle=\color{orange},
  numbers=left,
  numbersep=13pt,
  numberstyle=\small\color{mygray},
}
\lstset{escapechar=@,style=customc}
\linespread{1.45}
\newcommand{\HRule}{\rule{\linewidth}{0.5mm}}

\begin{document}

\begin{titlepage}
\begin{center}

\textsc{\Large Московский Физико-Технический Институт}\\
\textsc{\large (Национальный Исследовательский Университет)}\\[1.5cm]

% Upper part of the page. The '~' is needed because \\
% only works if a paragraph has started.
\includegraphics[width=0.25\textwidth]{img/logo.png}~\\[1cm]

\textsc{\Large Бакалаврская работа}\\[0.5cm]

% Title
\HRule \\[0.4cm]
{ \LARGE \bfseries Влияние вышестоящих регуляторных регионов и промоторов на связь между бактериофагами и их хозяевами \\[0.4cm] }

\HRule \\[1.5cm]

% Author and supervisor
\noindent
\begin{minipage}{0.4\textwidth}
\begin{flushleft} \large
\emph{Студент:}\\
Шарафутдинов Эмиль
\end{flushleft}
\end{minipage}%
\begin{minipage}{0.4\textwidth}
\begin{flushright} \large
\emph{Преподаватель:} \\
 Родригес Валера Франсиско Эдуардо
\end{flushright}
\end{minipage}

\vfill

% Bottom of the page
{\large 24 апреля 2021 г.}

\end{center}
\end{titlepage}

\newpage
\section{Аннотация}

Бактериофаги (фаги) - это вирусы, которые избирательно поражают бактериальные клетки. Такая специфика обусловлена
зависимостью фага от клеточного механизма хозяина, который нужен фагу для репликации и биосинтеза вирусных компонентов.
Пример такой зависимости  можно привести для фага Т4. Во время транскрипции, ранняя конкуренция за РНК-полимеразу хозяина
позволяет вирусу полностью захватить механизм репликации путем сигма-присвоения(sigma appropriation) [1]. Недавно
показали[2], что фаги, которые заражают цианобактерии, кодируют ген фотосинтеза. Причём транскрипция этого гена реагирует
на те же условия окружающей среды, что и ген, который кодируется у хозяина. Из этого можно предположить, что существует
связь между регуляторными элементами фага и регуляторными элементами их хозяина. Это поможет лучше предсказывать
взаимоотношения фаг-хозяин.

\newpage
\tableofcontents



\newpage
\section{Введение}

    В одной статье были собраны результаты разных исследований, касающихся взаимоотношений фага T4 и его хоста – E. Coli.
    В этих работах было показано, как ведёт себя фаг, который внедряется в бактерию.
	У фага Т4 нет собственных РНК-полимераз, поэтому он использует РНК-полимеразы хоста для транскрипции. Но почему РНК
	полимераза хоста должна сразу транскрибировать именно с ДНК Т4? Не должна. Но всё же обнаружили, что некоторая часть
	ДНК фага транскрибируется РНК-полимеразой хоста сразу. Эту часть транскрипции назвали ранней транскрипцией фага. На
	этом слайде можно видеть, почему эта ранняя транскрипция просходит. Здесь обозначены разные мотивы промоторов ДНК
	фага (посередине – ранняя, снизу - промежуточная) и хоста (сверху). Эти мотивы, которые находятся в промоторе,
	называются вышестоящими регуляторными элементами. Как видно, в промоторе ранней транскрипции регуляторные элементы
	между фагом и хостом совпадают, и выше -35 участка, что позволяет РНК полимеразе цепляться за ДНК фага также хорошо,
	как за ДНК хоста.  Всё, что до мотива TGn совпадает у всех промоторов. MotA Box в промоторе промежуточной
	транскрипции нужен для зацепки белка, который способствует промежуточной транскрипции. Вкратце, что там просходит: в
	этот момент с РНК, транскрибировавшейся от фага (в ранней транскрипции), транслируются белки, которые модифицируют
	РНК-полимеразы хозяина таким образом, что далее начинают транскрибироваться только ДНК фага. В том числе в начальной
	транскрипции ещё образуется белок MotA, который позволяет ДНК прикрепиться к модифицированной РНК-полимеразе. После
	чего с этой ДНК начинает транскрибироваться РНК. Этот этап уже называют промежуточной транскрипцией. Этот процесс
	присвоения РНК-полимеразы хозяина называют сигма присвоением. Ещё раз. Проиcходит это так:

	В самом начале транслируются белки, которые глушат РНК-полимеразу
	Вместе с этим транслируются белки, которые позволяют при заглушенной полимиразе транскрибироваться ДНК фага на этой
	полимеразе, они садятся на ДНК фага
	Таким образом фаг полностью захватывает РНК-полимеразу хозяина
	Рассмотри ещё один интересный случай для нас. Ещё в одной статье авторы смотрели, как ведут себя транкриптомы фагов и
	их хозяев (цианобактерий в этом случае) в зависимости от интенстивности света, которым они светили. При маленькой
	интенсивности транскриптомы делились в среднем поровну, но при бОльшей интенсивности среднее число РНК фагов было
	намного больше, чем у цианобактерий. Авторы решили посмотреть, что это за участок РНК, который так сильно возрастает
	у фага при сильном свете. Они обнаружили, что это участок РНК, ответственный за белок psbA, который оказался общим у
	фага и его хозяина. Авторы предположили, что это вполне возможно, если фаги позаимствовали гены у хозяев в ходе
	рекомбинации.

\newpage
\section{Цель и задачи} \label{sec:code}

Эти статьи – хорошие примеры, но они не единственные в своём роде. Нам стали интересны такие события, и в нашей работе
мы предположили, что такие схожие вышестоящие регуляторные участки в ДНК фагов и их хостов могут встречаться не только у
цианобактерий и их фагов и T4 с E. Coli, а у всех фагов и их хостов.
В том числе, мы посчитали, что у многих фагов и их хозяев могут встречаться общие белки. Таким образом, если мы сможем
установить взаимоотношения фагов и их хозяинов, мы в будущем сможем прогнозировать отношения фагов и каких-либо хостов.

\newpage
\section{Обзор литературы} \label{sec:math}
\subsection{Hinton D.M.}
    \subsubsection{Фон}
        \par{Контроль транскрипции имеет решающее значение для правильной экспрессии генов и упорядоченного развития. В течение
        многих лет бактериофаг Т4 обеспечивал простую модельную систему для исследования механизмов, регулирующих этот
        процесс. Развитие Т4 требует транскрипции ранних, средних и поздних РНК. Поскольку Т4 не кодирует свою собственную
        РНК-полимеразу, он должен перенаправить полимеразу своего хозяина, E.coli, в нужный класс генов в нужное время. Т4
        достигает этого за счет действия факторов, кодируемых фагами. Здесь я рассматриваю недавние исследования, изучающие
        транскрипцию пререпликативных генов T4, которые экспрессируются как ранние и средние транскрипты. Ранние РНК
        генерируются сразу после заражения из промоторов Т4, которые содержат отличные последовательности распознавания для
        полимеразы хозяина. Следовательно, ранние промоторы чрезвычайно хорошо конкурируют с промоторами-хозяевами за
        доступную полимеразу. Ранняя промоторная активность T4 дополнительно усиливается действием белка T4 Alt, компонента
        головки фага, который вводится в кишечную палочку вместе с ДНК фага. Alt модифицирует Arg265 на одной из двух
        субъединиц а РНК-полимеразы. Хотя работа с промоторами хозяина предсказывает, что эта модификация должна снизить
        активность промотора, транскрипция некоторых ранних промоторов Т4 увеличивается, когда РНК-полимераза модифицируется
        Алт. Транскрипция промежуточных генов Т4 начинается примерно через 1 минуту после заражения и протекает двумя путями:
        1) распространение ранних транскриптов в нижестоящие средние гены и 2) активация средних промоторов Т4 посредством
        процесса, называемого присвоением сигмы. При этой активации коактиватор Т4 AsiA связывается с регионом 4
        \(\sigma^{70}\), sсубъединицей специфичности РНК-полимеразы. Это связывание резко реконструирует эту часть
        \(\sigma^{70}\), что затем позволяет активатору Т4 MotA также взаимодействовать с \(\sigma^{70}\). Кроме того,
        перестройка \(\sigma^{70}\) предотвращает формирование нормальных контактов области 4 с областью -35 промоторной ДНК,
        что, в свою очередь, позволяет MotA взаимодействовать со свом сайтом связывания ДНК, коробкой MotA, центрированной в
        области -30 средней промоторной ДНК. Присвоение сигмы T4 показывает, как специфический домен в РНК-полимеразе может
        быть переформован, а затем использован для изменения специфичности промотора. \cite{hinton}}
        \par{
        Фоновая экспрессия генома Т4-это строго регулируемый и элегантный процесс, который начинается сразу после заражения
        хозяина. Основной контроль над этой экспрессией происходит на уровне транскрипции. Т4 не кодирует свою собственную
        РНК-полимеразу (РНАП), а вместо этого кодирует множество факторов, которые служат для изменения специфичности
        полимеразы по мере развития инфекции. Эти изменения коррелируют с временной регуляцией трех классов транскрипции:
        ранней, средней и поздней. Ранняя и средняя РНК обнаруживается дорепликативно
        \cite{hinton1,hinton2,hinton3,hinton4,hinton5,hinton6}, в то время как поздняя транскрипция сопровождается
        репликацией Т4 и обсуждается в другой главе. Ранние транскрипты T4 генерируются из ранних промоторов (Pe),которые
        активны сразу после заражения. Ранняя РНК обнаруживается даже в присутствии хлорамфеникола, антибиотика, который
        предотвращает синтез белка. Напротив,средние транскрипты Т4 генерируются примерно через 1 минуту после заражения при
        \(37^\circ C\)  и требуют синтеза фагового белка. Средняя РНК синтезируется двумя способами: 1) активация средних
        промоторов (Pm) и 2) расширение транскриптов Pe из ранних генов в нижестоящие средние гены. Этот обзор посвящен
        исследованиям ранней и средней транскрипции T4, начиная с тех, которые подробно описаны в последней книге T4
        \cite{hinton1}-\cite{hinton5}. На момент этой публикации ранние и средние транскрипты были подробно охарактеризованы,
        но механизмы, лежащие в основе их синтеза, только зарождались. В частности, эксперименты in vitro только что
        показали, что для активации средних промоторов требуется модифицированный Т4 РНАП и активатор Т4 МотА
        \cite{hinton7,hinton8}. Последующая работа определила необходимую модификацию RNAP как плотное связывание белка 10
        кДа, AsiA, с субъединицей \(\sigma^{70}\) RNAP \cite{hinton9}-\cite{hinton13}. Кроме того, в настоящее время стало
        доступно множество структурной и биохимической информации о RNAP E. coli [рассмотрено в [14-16]], MotA и AsiA
        [рассмотрено в \cite{hinton2}]. Как подробно описано ниже, теперь у нас есть гораздо более механистическое понимание
        процесса дорепликативной транскрипции Т4. Чтобы понять этот процесс, мы сначала начнем с обзора транскрипционного
        механизма хозяина и RNAP.} \\ 
    {\Large Механизм транскрипции кишечной палочки }
        \par{Голофермент E. coli RNAP, как и все бактериальные RNAP, состоит из ядра субъединиц
        \((\beta,\beta',\alpha_1,\alpha_2\) и \(\omega)\), которое содержит активный сайт для синтеза РНК, и фактора
        специфичности \(\sigma\), который распознает промоторы в ДНК и устанавливает начальный сайт для транскрипции.
        Первичный \(\sigma\), \(\sigma^{70}\) в E. coli, используется во время экспоненциального роста; альтернативные
        факторы \(\sigma\) направляют транскрипцию генов, необходимых в различных условиях роста или во время стресса
        [рассмотрено в [17-19]]. Анализ последовательности/функций сотен факторов \(\sigma\) выявил различные регионы и
        субрегионы сохранения. Большинство \(\sigma\)-факторов имеют сходство в областях 2-4, центральной через С-концевую
        часть белка, в то время как первичные \(\sigma\)-факторы также имеют родственную N-концевую часть, область 1.}
        \par{Недавняя структурная информация, наряду с предыдущими и текущими биохимическими и генетическими работами
        [рассмотренными в[14,15,20,21]], привела к биомолекулярному пониманию функции RNAP и процесса транскрипции. В
        настоящее время доступны структуры голофермента, ядра и частей первичных \(\sigma\) термофильных бактерий с ДНК и без
        ДНК [15,16,22-28], а также структуры областей только E. coli \(\sigma^{70}\) [29] и в комплексе с другими белками
        [26,30]. Эта работа показывает, что взаимодействие между \(\sigma^{70}\) и ядром внутри голофермента RNAP является
        обширным (рис. 1). Она включает в себя контакт между частью области \(\sigma\) 2 и спиральным/спиральным
        доменом,состоящим из \(\beta\), \(\beta'\), взаимодействие области \(\sigma^{70}\) 1.1 в“ челюстях” в нисходящем
        канале ДНК (где ДНК ниже по течению от места начала транскрипции будет расположена, когда RNAP связывает промотор), а
        также взаимодействие между областью \(\sigma^{70}\) 4 и частью субъединицы \(\beta\), называемой \(\beta\)-лоскутом.}
        \par{Для начала транскрипции части RNAP должны сначала распознать и связать с элементами распознавания двухцепочечной
        (ds) ДНК, присутствующими в промоторной ДНК(рис. 1) [рассмотрено в [20]]. Каждый из C-концевых доменов субъединиц
        \(\alpha\) (\(\alpha\)-CTDS) может взаимодействовать с элементом UP, богатыми последовательностями A/T,
        присутствующими между позициями -40 и -60. Части \(\sigma^{70}\), присутствующие в RNAP, могут взаимодействовать с
        тремя различными элементами dsDNA. Спираль-поворот-спираль, мотив связывания ДНК в области \(\sigma^{70}\) 4 может
        связываться с элементом -35, область \(\sigma^{70}\) 3 может связываться с последовательностью -15 TGn -13 (TGn), а
        субрегион \(\sigma^{70}\) 2.4 может связываться с позициями -12/-11 элемента -10. Распознавание элемента -35 также
        требует контакта между остатками в области 4 \(\sigma^{70}\) и заслонкой \(\beta\), чтобы правильно расположить
        \(\sigma^{70}\) для одновременного контакта элемента -35 и нижестоящих элементов. Как правило, промотор должен
        содержать только два из трех \(\sigma^{70}\)-зависимых элементов для активности; таким образом, E. промоторы coli
        могут быть свободно классифицированы как -35/-10 (основной класс), TGn/-10 (также называемый расширенным -10) или
        -35/TGn[рассмотрено в [20]].}
        \par{Первоначальное связывание RNAP с элементами промотора dsDNA обычно приводит к нестабильному, “закрытому”
        комплексу(RPc) (рис. 1). Создание стабильного, “открытого” комплекса(RPo) требует изгиба и разматывания ДНК [31] и
        основных конформационных изменений (изомеризации) полимеразы (рис.1) [32,33]; рассмотрено в[20]]. В RPo размотка ДНК
        создает пузырь транскрипции от -11 до ~+3, обнажая одноцепочечный(ss) шаблон ДНК для транскрипции. Добавление
        рибонуклеозидтрифосфатов (RNTPS) затем приводит к синтезу РНК, которая остается в виде гибрида ДНК/РНК в течение
        примерно 8-9 лет. Генерация более длинной РНК инициирует экструзию РНК через выходной канал РНК, образованный частями
        \(\beta\) и \(\beta'\) внутри ядра. Поскольку этот канал включает в себя связанный с \(\sigma^{70}\)
        \(\beta\)-лоскут, считается, что прохождение РНК через канал помогает высвобождению \(\beta\) из ядра, облегчая
        клиренс промотора. Полученный комплекс элонгации, EC, содержит основную полимеразу, шаблон ДНК и синтезированную РНК
        (рис. 1) [рассмотрено в [34]]. ЭК быстро перемещается по ДНК со скоростью около 50 нтл/сек, хотя комплекс может
        останавливаться в зависимости от последовательности [35]. Прекращение транскрипции происходит либо при внутреннем
        сигнале прекращения, структуре стволовой петли (шпильки), за которой следует U-богатая последовательность, либо при
        Rho-зависимом сигнале прекращения [рассмотрено в [36,37]]. Образование шпильки РНК внутренней последовательностью
        терминатора может способствовать прекращению путем дестабилизации гибрида РНК/ДНК. Rho-зависимая терминация
        опосредуется взаимодействием белка Rho с участком колеи (последовательностью использования Rho), неструктурированной,
        иногда богатой C последовательностью, которая находится выше по течению от места терминации. После связывания с РНК
        Rho использует гидролиз АТФ для транслокации вдоль РНК, догоняя EC в месте паузы. Как именно Ро разъединяет
        приостановленный комплекс, до конца не понятно; ДНК:РНК-геликазная активность Rho может обеспечить силу для
        “выталкивания” RNAP из ДНК. Только Rho достаточно для завершения в некоторых зависимых от Rho местах завершения.
        Однако на других участках процесс прекращения также нуждается в вспомогательных белках E. coli NusA и/или NusG
        [рассмотрено в [36]].}
        \par{При наличии в межгенных регионах участки колеи легко доступны для взаимодействия с Rho. Однако, когда они
        присутствуют в областях, кодирующих белки, эти сайты могут быть замаскированы путем трансляции рибосом. В этом случае
        прекращение Rho не наблюдается, если только вышестоящий ген не переведен, например, когда мутация породила
        нонсенс-кодон. В таком случае Rho-зависимая терминация может предотвратить распространение транскрипции в нисходящий
        ген. Таким образом, в этой ситуации, которая называется полярностью [38], предотвращается экспрессия как восходящего
        мутированного гена, так и нисходящего гена.}
    \subsubsection{Ранняя транскрипция Т4}
        {\Large Ранние промотеры}
            \par{Т4 заражает только экспоненциально растущую кишечную палочку, и транскрипция ранних генов Т4 начинается сразу
            после заражения. Таким образом, для эффективной инфекции фаг должен быстро перенаправить
            \(\sigma^{70}\)-ассоциированный RNAP, который активно участвует в транскрипции генома хозяина, на ранние
            промоторы T4. Это немедленное поглощение является успешным отчасти потому, что большинство ранних промоторов
            T4 содержат отличные соответствия элементам распознавания \(\sigma^{70}\)-RNAP (-35, TGn и -10 элементов) и
            элементам \(\alpha\)-CTD UP (рис.2; списки последовательностей ранних промоторов T4 см. в [4,5]). Однако
            выравнивание последовательностей ранних промоторов Т4 выявляет дополнительные области консенсуса,
            предполагая, что они содержат другие биты информации, которые могут оптимизировать взаимодействие RNAP
            хозяина с элементами промотора. Следовательно, в отличие от большинства промоторов хозяина, которые относятся
            к классу -35/-10, TGn/-10 или -35/TGn, ранние промоторы T4 могут быть описаны как “сверх” UP/-35/TGn/-10
            промоторов. Действительно, большинство ранних промоторов Т4 очень хорошо конкурируют с промоторами хозяина за
            доступный РНАП [39] и похожи на другие очень сильные промоторы фагов, такие как T7 \(P_{A1}\) и \(\lambda\)
            $ P_L $}\\
        {\LargeБелок Т4 Алт}
            \par{Помимо явной силы своих ранних промоторов, у T4 есть еще одна стратегия, белок Alt, для установления
            транскрипционного доминирования [[40-43], рассмотренная в [1,4]]. Алт, моно-АДФ-рибозилтрансфераза,
            АДФ-рибозилирует специфический остаток, Arg265, на одной из двух субъединиц \(\alpha\) РНАП. Кроме того, Алт
            модифицирует часть других белков-хозяев, включая другие субъединицы RNAP и белки-хозяева, участвующие в
            трансляции и метаболизме клеток. Алт-это внутренний белок головки фага, который вводится вместе с ДНК фага.
            Следовательно, модификация Алт происходит сразу после заражения и не требует синтеза фагового белка. Каждая
            субъединица \(\alpha\) различна (одна \(\alpha\) взаимодействует с \(\beta\), в то время как другая
            взаимодействует с \(\beta'\)), и считается, что модификация Alt специально нацелена на конкретный \(\alpha\),
            хотя какой именно \(\alpha\) неизвестен.}
            \par{Какова цель модификации Alt? Было показано, что основная цель Алт,\(\alpha\) Arg265, имеет решающее
            значение для взаимодействия \(\alpha\)-CTD с элементом промотора UP [44-46] и с некоторыми активаторами
            хозяина, включая белок рецептора c-AMP (CRP), глобальный регулятор E. coli [46,47]. Таким образом, очевидная
            гипотеза состоит в том, что Алт просто ухудшает работу промоторов хозяина, которые либо нуждаются в этих
            активаторах, либо усиливаются взаимодействием элементов \(\alpha\)-CTD/UP. Однако сверхэкспрессия Алт из
            плазмиды не влияет на Е. рост кишечной палочки [40], и общая транскрипция ДНК E. coli in vitro не нарушается
            при использовании модифицированного Alt RNAP [48]. Вместо этого, по-видимому, модификация Alt полезна,
            поскольку она повышает активность некоторых ранних промоторов Т4. Это 2-кратное повышение активности
            наблюдалось как in vivo [40,49], так и in vitro[48]. Как Alt-модификация стимулирует конкретные ранние
            промоторы, неизвестно, но ясно, что это связано не только с их общей силой. Другие сильные промоторы, такие
            как \(P_{tac}\), \(T7\) \(P_{A1}\) и \(P_{A2}\), T5 \(P_{207}\) и даже некоторые из ранних промоторов T4, не
            затрагиваются при использовании модифицированного Alt RNAP [49]. Алт-опосредованная стимуляция промотора
            также не зависит от специфических \(\sigma^{70}\)-зависимых элементов (-35, TGn и -10 элементов); некоторые
            промоторы с идентичными последовательностями в этих областях стимулируются Алт, в то время как другие-нет
            [49]. Всесторонний мутационный анализ ранних промоторов \(T4 P_{8.1}\) и \(P_{tac}\) показывает, что не
            существует единой специфической промоторной позиции(позиций), ответственной за эффект Алт. Этот результат
            предполагает, что механизм стимуляции Алт может включать перекрестное взаимодействие между RNAP и более чем
            одной промоторной областью [50] или что АДФ-рибозилирование $ \alpha $ Arg265 является вторичной, менее
            значимой активностью Алт и необходима дополнительная работа по важности этого введенного фермента.}\\
        {\Large Продолжение ранних стратегий доминирования Т4}
            \par{Поскольку промоторы Т4 настолько эффективны в конкурентной борьбе с хозяином, всплеск немедленной ранней
            транскрипции происходит в течение первой минуты заражения. Из этой транскрипции следует волна ранних
            продуктов, которые продолжают захват фагами транскрипционного механизма хозяина. Одним из таких продуктов
            является белок T4 Alc, терминатор транскрипции, специфичный для ДНК, содержащей dC, то есть ДНК, содержащей
            немодифицированные цитозины. Следовательно, Alc прекращает транскрипцию из ДНК хозяина, не влияя на
            транскрипцию из ДНК T4, цитозины которой гидроксиметилированы и глюкозилированы [[51,52]; рассмотрено в
            [1,4]]. Alc направляет RNAP на завершение в нескольких, частых и дискретных участках вдоль ДНК, содержащей
            dC. Механизм Alc неизвестен. В отличие от других завершающих факторов, Alc, по-видимому, не взаимодействует
            ни с РНК, ни с ДНК, и снижение скорости синтеза РНК или приостановка RNAP вблизи сайта завершения Alc
            фактически ухудшает завершение Alc [51]. Мутации в N-концевой области субъединицы $ \beta $ RNAP, области,
            которая не является существенной для E. coli (необязательная область I), предотвращают Alc -опосредованное
            прекращение, предполагая, что сайт взаимодействия для Alc может находиться в этой области [52].}
            \par{Т4 также кодирует два других АДФ-рибозилирующих фермента, ModA и ModB, в качестве ранних продуктов. 
            Как и Alt, ModA изменяет Arg265 RNAP $ \alpha $ [[53,48]; рассмотрено в [1,4]]. Однако, в отличие от Alt,
            ModA почти исключительно нацелена на субъединицы RNAP $\alpha$. Кроме того, ModA изменяет обе субъединицы
            $\alpha$, поэтому нет асимметрии в модификации ModA. Синтез ModA высокотоксичен для кишечной палочки. In
            vitro ModA-модифицированный RNAP не способен взаимодействовать с элементами UP или взаимодействовать с CRP
            [цитируется в [40]] и менее активен, чем немодифицированный RNAP при использовании ДНК E. coli или T4 [48].
            Таким образом, было высказано предположение, что ModA помогает уменьшить активность как хозяина, так и
            раннего промотора Т4, перепрограммируя транскрипционный механизм для предстоящей волны средней транскрипции
            [48]. Однако делеция гена modA не влияет на быстрое снижение ранней транскрипции или снижение синтеза ранних
            продуктов гена, которое начинается примерно через 3 минуты после заражения [54]. Этот результат говорит о
            том, что фаг использует другие, пока неизвестные стратегии, чтобы остановить транскрипцию от ранних
            промоторов. ModB, другой ранний АДФ-рибозилирующий фермент, нацелен на факторы трансляции хозяина,
            рибосомальный белок S30 и триггерный фактор, который, по-видимому, помогает стимулировать трансляцию Т4[43].}
            \par{Наконец, многие из ранних транскриптов включают гены неизвестной функции и происходят из областей генома
            Т4, которые не являются существенными для заражения кишечной палочкой дикого типа(wt) в нормальных
            лабораторных условиях. Предположительно, эти гены кодируют фаговые факторы, которые полезны в определенных
            условиях роста или у определенных штаммов. Неизвестно, помогает ли какой-либо из этих генных продуктов Т4 в
            его захвате транскрипционного механизма хозяина.} \\
        {\Large Переход на среднюю транскрипцию}
            \par{В течение минуты после заражения при $ 37^{\circ}C $некоторые из ранних продуктов Т4 опосредуют переход
            от ранней к средней экспрессии генов. Как подробно описано ниже, активатор MotA и соактиватор AsiA являются
            важными партнерами в этом переходе, поскольку они направляют RNAP для транскрипции от средних промоторов.
            Кроме того, белок CoMC-$\alpha$, описанный ниже, также может играть роль в распространении ранних РНК в
            нижестоящие средние гены или стабильности таких транскриптов после их формирования.}
            \par{Когда начинается средняя транскрипция, некоторые ранние РНК быстро распадаются после их первоначального
            всплеска транскрипции. Это происходит из-за активности раннего генного продукта RegB, эндорибонуклеазы,
            которая специфически нацелена на некоторые ранние мРНК Т4. Для мРНК MotA и самого RegB сайт расщепления RegB
            находится в последовательности Шайна-Далгарно; для мРНК CoMC-$\alpha$ сайт находится в богатых AU
            последовательностях выше и ниже этой последовательности [55]. Механизм, с помощью которого RegB распознает и
            выбирает конкретный участок расщепления, пока неизвестен.}
            \par{Начало средней транскрипции T4 также завершает процесс устранения транскрипции хозяина, просто удаляя
            шаблон ДНК хозяина для RNAP. Т4-кодируемые нуклеазы, в первую очередь ЭндоИИ, кодируемые DeNA, и ЭндоИВ,
            кодируемые DeNB, избирательно разрушают ДНК хозяина, содержащую dC ([56,57] и ссылки в них). Таким образом,
            через несколько минут после заражения, по существу, нет ДНК хозяина, которую можно было бы транскрибировать.}
    \subsubsection{Транскрипция средних генов из средних промоторов Т4}
        {\Large Средние промотеры}
        \par{Средние гены в основном кодируют белки, необходимые для репликации, рекомбинации и метаболизма нуклеотидов;
        различные ТРНК, кодируемые Т4; и факторы транскрипции, которые программируют переход от средней к поздней
        активации промотора. Средние РНК возникают по 2 путям: расширение ранней транскрипции в средние гены (обсуждается
        позже) и активация средних промоторов Т4 с помощью процесса, называемого $\sigma$-присвоением [2]) На сегодняшний
        день было идентифицировано почти 60 средних промоторов (таблица 1). В отличие от ранних промоторов, средние
        промоторы T4 содержат элемент хозяина, $\sigma^{70}$-зависимую последовательность -10 и элемент фага, блок MotA,
        который центрирован на -30 и заменяет $\sigma^{70}$-зависимый элемент -35, присутствующий в ранних промоторах T4
        и большинстве промоторов хозяина (рис.2). Кроме того, около половины средних промоторов также содержат TGn,
        расширенную последовательность -10. Активация средних промоторов фага требует согласованных усилий двух ранних
        продуктов Т4, AsiA и MotA.}

\subsection{Energy limitation of cyanophage development: implications formarine carbon cycling}
    \subsubsection{Введение}
        \par{Понимание реакции биосферы на изменение климата требует детального знания о биологической трансформации
        углерода на Земле. Стало очевидным, что океан представляет собой важный поглотитель атмосферного углекислого газа
        ($CO_2$)[1]. Только два рода пикофитопланктона доминируют в режимах открытого океана-Prochlorococcus и
        Synechococcus[2-4]. В целом, эти близкородственные роды ответственны за $ ~25\% $ фиксации океанических $CO_2$
        [2]. Эти организмы стали моделью для изучения потока углерода из $CO_2$ в микробную петлю, а также
        функционирования планктонных морских сообществ в целом [5,6]. Однако показатели биотических потерь
        (антагонистические взаимодействия, выпас скота, вирусный лизис) прохлорококка и синехокка остаются
        малоизученными. В случае вирусов эта точка зрения дополняется открытием генов, кодируемых вирусами, которые могут
        действовать для поддержания фотосинтеза во время инфекции [7,8], так что, несмотря на окончательную потерю
        фиксированного углерода в растворенном органическом веществе в результате лизиса, фиксация $CO_2$ может
        поддерживаться временно в течение относительно длительных латентных периодов вируса. Недавно было показано, что
        на самом деле цианофаг отключает $CO_2$ инфекцию на ранних стадиях заражения, но при этом поддерживает
        фотосинтетические световые реакции [9]. Преобладает мнение, что цианофаги используют избыток АТФ и восстановителя
        от фотофосфорилирования для подпитки репликации ДНК, одновременно ингибируя дорогостоящий цикл Кальвина Бенсона
        [10,11]. Поэтому мы предположили, что репликация ДНК является ограничивающим фактором морфогенеза цианофагов и
        что увеличение скорости фотофосфорилирования в результате увеличения облучения резко изменит кинетику репликации
        ДНК фага и, как следствие, продуктивность фага. Такие изменения будут иметь важные последствия для нашего
        понимания вызванного вирусами давления лизиса сообществ открытого океана, учитывая широкий диапазон светового
        потока in situ. \cite{puxty-evanx}}
    \subsubsection{Результаты}
        {\Large Модули интенсивности света кинетика инфекции и фотофизиология}
            \par{Мы предположили, что перенос инфицированных клеток синехококка на высокий свет (HL) обеспечит
            повышенную энергию для репликации фагов, и поэтому мы будем наблюдать увеличение скорости репликации ДНК
            цианофагов. Чтобы проверить эту идею, мы провели эксперименты со сдвигом света с помощью синехококка.
            WH7803, инфицированный цианофагом S-PM2d [12], и контролировал развитие фага. Мы провели анализ qPCR,
            количественно определяя копии хромосомы фага из внутриклеточных и внеклеточных фракций, чтобы измерить
            скорость синтеза ДНК внутри клеток и определить латентные периоды и размеры пакетов}
            \par{Мы наблюдали значительную разницу в развитии цианофагов в HL по сравнению с LL (рис.1). При ГЛ
            латентный период был сокращен на 5 ч (уменьшение длины на 40\%), о чем свидетельствует как более раннее
            уменьшение копий внутриклеточной ДНК (рис.1а), так и более раннее увеличение числа копий внеклеточной ДНК
            (рис.1б). Этот ранний всплеск в условиях HL также был очевиден из мониторинга мутности культуры (рис.1c).
            Для сравнения, не было никакой существенной разницы в размере всплеска между различными световыми
            обработками (рис.1b). Удивительно, но, вопреки нашей гипотезе, не было никакой разницы во времени, длине
            или скорости синтеза ДНК фага между этими двумя условиями (рис.1а). Мы постулировали, что увеличение
            фотохимически управляемого синтеза АТФ в условиях ГЛ обеспечит больше энергии и ресурсов для подпитки
            репликации ДНК фага. Действительно, эта идея была задействована [13] и явно смоделирована [14,15] как
            объяснение приобретения цианофагами генов фотосинтеза, особенно учитывая, что другие стрессы, такие как
            ограничение фосфатов, оказывают заметное влияние на синтез ДНК цианофагов [16]. Наша гипотеза основана на
            предположении, что фотохимия поддерживается таким образом, что избыток поглощенных фотонов может быть
            использован для создания трансмембранного потенциала для синтеза АТФ. Поэтому мы проверили это, измерив
            максимальный квантовый выход фотосистемы II(PSII) $ (F_v/F_m)$ во время эксперимента. Этот анализ измеряет
            различия в "флуоресцентном транзиторе" хлорофилла от PSII. По мере того как реакционные центры,
            адаптированные к темноте, "закрываются" из-за уменьшения первичного акцептора электронов, происходит
            сопутствующее увеличение флуоресценции хлорофилла. Поэтому параметр $ F_v/F_m $ пропорционален числу
            открытых реакционных центров в данный момент времени [17]. Наши измерения фотохимического выхода PSII не
            показали никаких изменений в условиях LL в инфицированных или неинфицированных клетках (рис.2а). Для
            сравнения, наблюдалось быстрое снижение $F_v/F_m$ в неинфицированных клетках при HL, свидетельствующее о
            фотоповреждении реакционного центра PSII [18]. Через 8 ч при HL $F_v/F_m$ в неинфицированных клетках
            стабилизируется примерно на 65\% от уровня культур LL.Аналогично, в условиях HL наблюдалось начальное
            снижение $ F_v/F_m $ в инфицированных клетках с той же скоростью, что и в неинфицированных клетках
            (рис.2а). Однако затем наблюдалось увеличение $ F_v/F_m $ через 2 часа после заражения до ~90\% от уровня в
            LL, где он оставался до конца эксперимента. Это наблюдение говорит о том, что не только реакционные центры
            PSII активны в клетках, инфицированных HL, но и либо скорость их повреждения снижается, либо оборот PSII
            намного больше, чтобы поддерживать стабильную фотохимию при этом облучении.}
            \par{Поддержание фотохимии PSII при повышенном облучении предполагает, что в этих клетках было доступно
            больше энергии за счет фотофосфорилирования [19]. Тем не менее, несмотря на это, мы не наблюдали заметных
            изменений в скорости, времени или абсолютном количестве синтеза ДНК фага (рис.1а, б). Однако мы наблюдали
            сокращение латентного периода, и поэтому мы предполагаем, что может быть некоторый временной контроль
            лизиса в ответ на свет. Действительно, было показано, что бактериофаги обладают сложными механизмами
            задержки лизиса в ответ на потенциальные стрессы, возникающие во время инфекционного процесса [20].}\\
        {\Large Транскрипционный ответ цианофага на высокие условия освещения}
            \par{Чтобы расшифровать потенциальные регуляторные механизмы, возникающие на уровне транскрипции в условиях
            HL по сравнению с условиями LL, мы провели РНК-Seq-анализ транскриптомов из того же эксперимента. РНК
            экстрагировали из трипликатных культур через 1, 3, 6 и 9 ч после заражения в инфицированных клетках HL и
            LL, а образцы из каждой световой обработки объединяли в эквимолярных количествах для оценки общих изменений
            экспрессии. Считывания были сопоставлены как с S-PM2d, так и с Synechococcussp. Оценены эталонные геномы
            WH7803 и среднее значение на базовое покрытие. Наблюдалось значительное увеличение охвата транскриптома
            фага в условиях HL по сравнению с условиями LL (увеличение в 2,34 раза,$t_3=3,98$, p=0,014), что не было
            заметно у хозяина (рис.2b). Это было несмотря на то,что доля клеток, подвергшихся лизису в течение 9 ч в
            HL, была меньше по сравнению с отсутствием лизиса в условиях LL (рис.1а, б). Это свидетельствует о
            генерализованном увеличении скорости транскрипции фагов при ГЛ, несмотря на отсутствие изменений в
            репликации ДНК. Мы использовали дифференциальные тесты экспрессии, чтобы найти светочувствительные гены в
            геноме цианофага. К нашему удивлению, только один ген цианофага был пропорционально дифференцированно
            экспрессирован в условиях HL по сравнению с LL. Это был фотосинтетический AMG psbA, кодирующий полипептид
            D1, расположенный в ядре реакционного центра PSII (рис.3а)}
            \par{Кодируемый цианофагом psbA является " почти основным’ геном в цианомовирусах [21], результатом
            повторных событий приобретения от их непосредственных хозяев [22]. Таким образом, возможно, что
            приобретение фагами генов хозяина путем рекомбинации также кооптировало те же светозависимые регуляторные
            особенности psbA. Действительно, было показано, что другие цианофаги регулируют экспрессию АМГ, поглощающих
            фосфор, в ответ на ограничение фосфата, хотя и используют собственную двухкомпонентную систему PhoBR
            хозяина [16]. Цианобактерии, включая морской синехококк, демонстрируют светозависимую транскрипцию PSBA
            [18, 23]. В модельных цианобактериях несколько факторов влияют на светозависимое накопление транскриптов
            psbA, включая альтернативные сигма-факторы [24,25,26,27,28] и третичную структуру ДНК [29]. Более того,
            продукты деградации D1 непосредственно связываются перед последовательностями psbA, так что индуцированное
            светом повреждение может непосредственно влиять на транскрипцию psbA с положительной обратной связью [30].
            Мы не смогли идентифицировать какие-либо консервативные мотивы в восходящей области psbA S-PM2d, которые
            разделяются с их хозяевами-синехокками, возможно, предполагая различные стратегии регуляции между версиями
            хозяина и фага.}
            \par{Кроме того, S-PM2d имеет нехарактерно длинную межгенную последовательность выше psbA (232 п. н., по
            сравнению с медианой 6 п. н. по всему геному; рис. 3b). Это различие также верно для других цианомовирусов,
            где области psbA выше по течению варьируют в пределах 125-453 п. п. Поэтому мы также стремились обнаружить
            мотивы последовательности ДНК, общие для цианомовирусов, которые могут быть вовлечены в регуляцию psbA. Мы
            обнаружили шесть мотивов с высокой степенью достоверности (рис. S1), из которых два присутствовали в S-PM2d
            (рис. 3c). Эти два мотива содержат -35 (мотив 6) и -10 (мотив 3) элементов сайтов связывания фактора
            транскрипции $\sigma^{70}$, типичных для Т4-подобных ранних фаговых генов [31]. Кроме того, мотив 3 содержит сайт
            связывания для T4-подобного позднего гена сигма-фактора Gp55 [32] (рис. 3c). Таким образом, S-PM2d
            сохраняет мотивы, необходимые для скоординированной экспрессии psbA при ранней и поздней инфекции.}
            \par{Однако область psbA S-PM2d содержит новую 3' - антисмысловую нкРНК, cfrI, которая, как известно,
            экспрессируется [33], а также ранее не распознанный инвертированный повтор 36 bp, производящий вероятную
            петлю шпильки РНК в 5' - нетранслируемой области (5'-UTR; рис. 3d). Хорошо известно, что структурированные
            5'-UTR влияют на стабильность мРНК бактерий и бактериофагов [34], как и цис-кодируемые антисмысловые РНК
            [35], особенно для цианобактериальных PSBA [36]. Таким образом, эти особенности могут способствовать
            регуляции psbA цианофага в ответ на свет.}
            \par{В еще одном повороте S-PM2d psbA кодирует интрон самосплайсинга группы I в своем транскрипте psbA,
            особенность, которая, как оказалось, широко распространена в генах psbA цианофага из морских метагеномных
            библиотек [33, 37]. Интрон прерывает последовательность распознавания нисходящей самонаводящейся
            эндонуклеазы F-CPhI, тем самым обеспечивая само иммунитет, позволяя эндонуклеазе распространяться среди
            аллелей без интрона [38]. Интрон содержит несколько внутрикадровых стоп-кодонов. Таким образом,
            маловероятно, что несвязанный транскрипт кодирует функциональный полипептид D1 [33]. Тем не менее, интрон
            эффективно сращивается во время инфекции [33].}
            \par{Чтобы понять, играет ли сплайсинг интронов роль в экспрессии psbA, зависящей от света, и действительно
            проверить наши данные по РНК-Seq, мы провели количественную ПЦР с обратной транскриптазой (RT-qPCR) мРНК
            psbA цианофага. Мы разработали анализ сплайсинга для определения относительной экспрессии обеих изоформ
            транскрипта (рис. 4а). Мы использовали 5'-нуклеазный анализ qPCR, который дает специфичность каждой
            изоформе транскрипта с использованием зонда, связанного с флюорогасителем 25-30 бп, который либо является
            внутренним по отношению к интрону (обнаружение несвязанной изоформы), либо соединяет соединение 1/2
            экзона(обнаружение сплайсированной изоформы). Обратный праймер также расположен либо в интроне (обнаружение
            несвязанной изоформы), либо в экзоне 2 (обнаружение сплайсированной изоформы) и специфичен для
            клонированных копий каждой изоформы.}
            \par{данные HE RT–qPCR подтверждают данные, полученные из РНК-Seq, показывая повышенную экспрессию как
            сплайсированных, так и несвязанных изоформ в HL (рис. 4b, c). Это увеличение особенно выражено в
            сплайсированной изоформе во время поздней инфекции, где мы наблюдали 10,6-кратное увеличение HL по
            сравнению с LL (рис. 4с), по сравнению только с 3,37-кратным увеличением несращенной изоформы. Важно
            отметить, что наш анализ позволяет определить соотношение сплайсированной и несвязанной изоформы (рис. 4d),
            которое показывает, что на протяжении ранней инфекции до 4 ч фаг поддерживает увеличение количества
            сплайсированной изоформы в ~2 раза (рис.4d). 4d). При ГЛ это сохраняется на протяжении всей инфекции, в то
            время как при ЛЛ мы наблюдали быстрое снижение примерно до эквимолярных величин к 6 часам. Таким образом,
            сплайсинг интрона, по-видимому, реагирует на интенсивность света. Насколько нам известно, это первое
            сообщение о дифференцированно сплайсированном интроне у бактериофага, хотя интересно, что эукариотическая
            водоросль Chlamydomonas reinhardtii содержит четыре интрона группы I в своем гене psbA [39], которые, как
            известно, дифференцированно сплайсируются в ответ на свет [40]. Поскольку несвязанные транскрипты,
            вероятно, не кодируют функциональные полипептиды, мы предполагаем, что дифференциальный сплайсинг мРНК psbA
            представляет собой механизм регуляции генов, потенциально в координации с антисмысловой нкРНК cfrI и/или
            структурированной 5'-UTR, влияющей на стабильность мРНК.}\\
            {\Large Дифференциальная экспрессия генов хозяина усиливает узкие места в трансляции}
            \par{В дополнение к дифференциальной экспрессии генов цианофагов мы также проанализировали светочувствительные гены хозяина синехококка. Из 39 генов, активированных в условиях высокой освещенности, 25 (64\%) были вовлечены в трансляцию или посттрансляционную модификацию, оборот белка и шапероны (таблица 1). Это включало 13/31 и 4/22 белков, образующих 50-ю и 30-ю субъединицы рибосомы, соответственно. Кроме того, обе субъединицы комплекса GroEL/GroES, а также шапероны DnaJ и DnaK были повышены (таблица 1). Между тем, три из четырех копий гена psbA хозяина также были активированы при высоком освещении (таблица 1), возможно, дополняя фаговые поставки мРНК psbA. Пониженные гены в основном включали компоненты аппарата сбора света, потенциальные гены подвижности клеток и субъединицы, участвующие в линейном транспорте электронов, в первую очередь субъединицы реакционного центра фотосистемы I (таблица 1).}
\newpage
\section{Материалы и Методы} \label{sec:code}
\subsection{Программные средства}
\begin{itemize}
    \item Virus-Host DB - База данных Virus-Host организует данные о взаимоотношениях между вирусами и их хостами,
    представленные в виде пар идентификаторов таксономии NCBI для вирусов и их хостов. База данных Virus-Host охватывает
    вирусы с полными геномами, хранящимися в 1) NCBI/RefSeq и 2) GenBank, номера присоединения которых указаны в геномах
    EBI. Информация о хосте собирается из RefSeq, GenBank (в свободном текстовом формате), UniProt, ViralZone и вручную
    курируется с дополнительной информацией, полученной в результате опросов литературы. \cite{virus-host}
    
    \item GeneMarkS-2 - программное обеспечение для поиска генов внутри геномов. \cite{lomsad}
    
    \item NCBI Entrez Programming Utilities - программное обеспечение, набор из восьми серверных программ, которые
    обеспечивают стабильный интерфейс с системой запросов и баз данных Entrez в Национальном центре биотехнологической
    информации (NCBI) \cite{entrez}
    
    \item Biopython - набор свободно доступных инструментов для биологических вычислений, написанных на Python
    международной командой разработчиков. \cite{biopython}
    
    \item The MEME Suite - программное обеспечение для поиска новых, нераскрытых мотивов (повторяющихся паттернов
    фиксированной длины) в последовательностях. МЕМ разбивает паттерны переменной длины на два или более отдельных
    мотива. \cite{bailey}
    
    \item Python 3.9.4 - среда для запуска скриптов формата *.py \cite{python}
    
    \item Jupyter Notebook -веб-приложение с открытым исходным кодом, которое позволяет создавать и обмениваться
    документами, содержащими живой код, уравнения, визуализации и повествовательный текст. Использование включает в себя:
    очистку и преобразование данных, численное моделирование, статистическое моделирование, визуализацию данных, машинное
    обучение и многое другое. \cite{jupyter}
    
    \item Скрипты отсюда \(https://github.com/poolsar42/phages-and-hosts\) - подумать, что про них написать!
    \cite{github}
    
    \item \(Fasta\_merge.ipynb\) - скрипт для объединения регуляторных областей хозяина и его фагов в один файл.
    \item \(Hosts\_upstream\_regions.ipynb\) - скрипт для выделения регуляторных областей.
    \item \(Proteins\_formation.ipynb\) - скрипт для сбора белковых последовтельностей в один файл.
    \item \(Searching\_host\_genomes.ipynb\) - скрипт для выгрузки необходимых к работе геномов хостов.
    \item \(Unpacking.ipynb\) - скрипт для распаковки выгруженных геномов.
    \item \(VH\_tsv\_writer.ipynb\) - скрипт, выписывающий в файл ID хоста и ID прилежащих к нему фагов.
    \item \(Virus\_genes.ipynb\) - скрипт для выгрузки геномов фагов.
    \item \(missing\_phage\_genomes.ipynb\) - скрипт для поиска фагов, для которых GMS2 не находит гены.
    \item \(proteins\_convertation.ipynb \) - скрипт для преобразования \\ ДНК-последовательностей
                                в белковые.
    \item \(regex\_promoter.ipynb \) - скрипт для анализа полученных результатов
    \item \(run\_gms2.py\) - скрипт для поиска генов в геномах
    \item \(tsv\_write.ipynb\) - скрипт для выписывания информации о связи между хостами и фагами в один файл
\end{itemize}

\subsection{Получение общих мотивов}
    \subsubsection{Протокол выгрузки геномов}
    \begin{itemize}
        \item С FTP сервера Virus-Hist DataBase выгрузить файл virushostdb.tsv - в нём содержится информация о связи между фагами и их хостами.
    
        \item Запустить скрипт \(tsv\_write.ipynb\) - выписывает в удобном формате все взаимоотношения между фагами и хостами.
    
      \item Запустить скрипт \(Virus\_genes.ipynb\) - для выгрузки геномов фагов.
        \item Запустить скрипт \(Searching\_host\_genomes.ipynb\) - для выгразки геномов хостов.
        \item Запустить скрипт \(Unpacking.ipynb\).

    \end{itemize}

    \subsubsection{Протокол получение вышестоящих регуляторных регионов}
        \begin{itemize}
            \item Запустить скрипт \(run\_gms2.py\)
            \item Запустить скрипт \(Hosts\_upstream\_regions.ipynb\)
        \end{itemize}
    \subsubsection{Протокол подготовки данных к поиску мотивов}
        \begin{itemize}
            \item Запустить скрипт \(Fasta\_merge.ipynb\)
            \item Запустить MEME Suite с получившимися файлами.
        \end{itemize}

\par{Нам нужно было посмотреть на сходтсва между вышестоящими регуляторными участками у хостов и фагов. Для этого нужно
было посмотреть на геномы хостов и фагов. Мы собрали геномы всех фагов, которые находятся в базе данных Virus-Host DB,
там для каждого фага также был указан или были указаны хосты, к которым он может прикрепиться. Их RefSeq ID. Мы всё
выгрузили с FTP сервера.}

\par{Далее мы выгрузили все геномы хостов с базы данных RefSeq NCBI. Когда у нас были геномы фагов и их хостов, нам было
нужно получить с этих геномов гены. Гены мы находили с помощью Gene-Marks 2. Что хорошо, Gene-Marks 2 также дал позицию
каждого гена в геноме. Каждому гену посчитали его транслирующийся белок, также с помощью Gene-Marks 2. После этого для
каждого гена нашли вышестоящие регуляторные области, длиной в 50 нуклеотидов, которые нужны для посадки РНК-полимераз(?).
Наша гипотеза как раз заключается в том, что эти вышестоящие области у фага и его хозяина для некоторых генов будут
совпадать по некоторым мотивам.}

\par{Мы поработали с последовательностями. Мы выгружали геномы всех фагов и хостов отдельными файлами. Мы предварительно
обработали файлы, достали из них вышестоящие регуляторные области длинной в 50 нуклеотов. Вышестоящие области хоста и
всех фагов, прилегающих к нему, мы закинули в один файл - в этом файле мы и искали общие мотивы.}

\par{Далее проверку на совпадения вышестоящих областей мы засунули все вышестоящие области хоста и его фагов в программу
MEME Suite. Как и ожидалось, мы получили множество профилей, которые были общими у фагов и их хоста. }

\par{Также MEME-Suite для каждого профиля мотива дала нам PFM (Position Frequency Matrix), поэтому у нас была возможность
сравнить сами мотивы между собой. У нас появилась гипотеза, что мотивы фагов, атакующих одного хоста, могут быть также
схожи и как-то коррелировать между собой. Мы решили посмотреть корреляцию 2D массивов - потому что для каждой позиции в
матрице частот у нас было 4 значения - на 4 возможных нуклеотида.}

\newpage
\section{Результаты и Обсуждение}

\newpage
\section{Выводы}


\newpage
\begin{thebibliography}{9}
    \addcontentsline{toc}{section}{\refname}
    \bibitem{hinton1} Stitt B, Hinton DM: Regulation of middle-mode transcription.In Molecular biology  of  bacteriophage  T4.Edited by: Karam JD, Drake J, Kreuzer KN, Mosig G, Hall D, Eiserling F, Black L, Spicer E, Kutter E, Carlson K, Miller ES.Washington, D.C.: American Society for Microbiology; 1994:142-160
    \bibitem{hinton2} Brody E, Rabussay D, Hall D:Regulation of transcription of prereplicative genes. In Bacteriophage  T4.Edited by: Mathews CK, Kutter EM, Mosig G,Berget PB. Washington, D. C.: American Society for Microbiology;1983:174-183
    \bibitem{hinton3} Hinton DM, Pande S, Wais N, Johnson XB, Vuthoori M, Makela A, Hook-Barnard I: Transcriptional takeover by sigma appropriation: remodelling of the sigma70 subunit of Escherichia coli RNA polymerase by the bacteriophage T4 activator MotA and co-activator AsiA.Microbiology2005,151:1729-1740.
    \bibitem{hinton4} Miller ES, Kutter E, Mosig G, Arisaka F, Kunisawa T, Ruger W: BacteriophageT4 genome. Microbiol  Mol  Biol  Rev2003,67:86-156
    \bibitem{hinton5} Wilkens K, Ruger W: Transcription from early promoters. In Molecular Biology  of  Bacteriophage  T4.Edited by: Karam JD, Drake JW, Kreuzer KN,Mosig G, Hall DH, Eiserling FA, Black LW, Spicer EK, Kutter E, Carlson K,Miller ES. Washington, D. C.: American Society for Microbiology;1994:132-141.
    \bibitem{hinton6} Weisberg R, Hinton DM, Adhya S: Transcriptional Regulation in Bacteriophage. In Encyclopedia  of  Virology.3 edition. Edited by: Mahy BWJ,van Regenmortel MHV. Oxford: Elsevier; 2008:174-186, 174-186.
    \bibitem{hinton7} Hinton DM: Transcription from a bacteriophage T4 middle promoter using T4 motA protein and phage-modified RNA polymerase.J  BiolChem 1991, 266:18034-18044
    \bibitem{hinton8} Schmidt RP, Kreuzer KN: Purified MotA protein binds the -30 region of a bacteriophage T4 middle-mode promoter and activates transcription invitro. J  Biol  Chem1992, 267:11399-11407
    \bibitem{hinton9} Stevens A: New small polypeptides associated with DNA-dependent RNA polymerase of Escherichia coli after infection with bacteriophage T4. Proc  Natl  Acad  Sci  USA1972, 69:603-607.
    \bibitem{hinton10} Stevens A:An inhibitor of host sigma-stimulated core enzyme activitythat purifies with DNA-dependent RNA polymerase of E. coli followingT4 phage infection.Biochem  Biophys  Res  Commun1973,54:488-493.
    \bibitem{hinton11} Ouhammouch M, Orsini G, Brody EN: The asiA gene product of bacteriophage T4 is required for middle mode RNA synthesis.J  Bacteriol1994, 176:3956-3965
    \bibitem{hinton12} Ouhammouch M, Adelman K, Harvey SR, Orsini G, Brody EN: BacteriophageT4 MotA and AsiA proteins suffice to direct Escherichia coli RNA polymerase to initiate transcription at T4 middle promoters. Proc  NatlAcad  Sci  USA1995, 92:1451-1455.
    \bibitem{hinton13} Hinton DM, March-Amegadzie R, Gerber JS, Sharma M: Bacteriophage T4middle transcription system: T4-modified RNA polymerase; AsiA, a sigma70 binding protein; and transcriptional activator MotA.Methods  Enzymol1996, 274:43-57
    \bibitem{virus-host} https://www.genome.jp/virushostdb/ 
    \bibitem{biopython} https://biopython.org/
    \bibitem{python} https://docs.python.org/3/
    \bibitem{jupyter} https://jupyter.org/
    \bibitem{github} https://github.com/poolsar42/phages-and-hosts
    \bibitem{entrez} https://www.ncbi.nlm.nih.gov/books/NBK25501/
    \bibitem{hinton} Hinton D.M. Transcriptional control in the prereplicative phase of T4 development // Virology
    journal. 2010. V. 7. P. 289. https://doi.org/10.1186/1743-422X-7-289
    \bibitem{puxty-evanx} Puxty, R.J., Evans, D.J., Millard, A.D. et al. Energy limitation of cyanophage development:
    implications for marine carbon cycling // ISME J. 2018, 12, 1273–1286. https://doi.org/10.1038/s41396-017-0043-3
    \bibitem{lomsad} Lomsadze, A., Gemayel, K., Tang, S., and Borodovsky, M. Modeling leaderless transcription and
    atypical genes results in more accurate gene prediction in prokaryotes // Genome research,  2018. 28(7), 1079–1089.
    https://doi.org/10.1101/gr.230615.117
    \bibitem{bailey} Bailey, T.L., Johnson, J., Grant, C.E., and Noble, W.S. The MEME Suite // Nucleic acids research,
    2015, 43(W1), W39–W49. https://doi.org/10.1093/nar/gkv416
    \bibitem{ryasik} Ryasik, A., Orlov, M., Zykova, E., Ermak, T., Sorokin, A. Bacterial promoter prediction: Selection
    of dynamic and static physical properties of DNA for reliable sequence classification // 2018
    https://www.worldscientific.com/doi/abs/10.1142/S0219720018400036
    
\end{thebibliography}

\end{document}

\usepackage[english,russian]{babel}